\documentclass[11pt,letterpaper]{article}

\usepackage[T1]{fontenc}
\usepackage[utf8]{inputenc}
\usepackage{mathpazo}
\usepackage[margin=1.15in]{geometry}
\usepackage[expansion=false]{microtype}
\usepackage{amsmath,amssymb}
\usepackage{booktabs}
\usepackage{enumitem}
\usepackage{titlesec}
\usepackage[dvipsnames,table]{xcolor}
\usepackage{tcolorbox}
\usepackage{fancyhdr}
\usepackage[hidelinks]{hyperref}
\tcbuselibrary{skins,breakable}

\definecolor{accent}{HTML}{1F4E79}
\definecolor{warn}{HTML}{A03A2B}
\definecolor{soft}{HTML}{F2F4F7}
\definecolor{rule}{HTML}{9AA5B1}

\linespread{1.05}
\setlength{\parskip}{0.35em}
\setlength{\parindent}{0pt}

\titleformat{\section}{\Large\bfseries\color{accent}}{\thesection}{0.7em}{}
\titleformat{\subsection}{\large\bfseries}{\thesubsection}{0.6em}{}
\titleformat{\subsubsection}{\normalsize\bfseries\itshape}{}{0em}{}
\titlespacing*{\section}{0pt}{2.2em}{0.8em}

\pagestyle{fancy}
\fancyhf{}
\fancyhead[L]{\footnotesize\color{rule}\textsc{Finding Mind}, draft 9 --- The Turn}
\fancyhead[R]{\footnotesize\color{rule}\thepage}
\renewcommand{\headrulewidth}{0.3pt}

\newtcolorbox{voice}[1]{
  breakable, enhanced, colback=soft, colframe=accent,
  boxrule=0pt, leftrule=2.5pt, arc=0pt, outer arc=0pt,
  left=10pt, right=10pt, top=8pt, bottom=8pt,
  title={\bfseries\color{accent}#1}, coltitle=accent,
  fonttitle=\bfseries, colbacktitle=soft, titlerule=0pt,
  attach boxed title to top left={xshift=10pt, yshift=-2pt},
  boxed title style={colback=soft, boxrule=0pt}
}

\newtcolorbox{flag}{
  breakable, enhanced, colback=white, colframe=warn,
  boxrule=0.9pt, arc=1pt, left=9pt, right=9pt, top=7pt, bottom=7pt
}

\newcommand{\loc}[1]{{\small\ttfamily\color{accent}\hyphenchar\font=`\-\relax #1}}
\setlength{\emergencystretch}{3.5em}
\newcommand{\fnd}[1]{\textbf{\color{warn}#1}}
\newcommand{\AI}{\mathrm{AI}}
\newcommand{\CN}{\mathrm{CN}}
\newcommand{\Semi}{\mathbb{S}}
\providecommand{\llbracket}{[\![}
\providecommand{\rrbracket}{]\!]}

\begin{document}

\begin{center}
{\LARGE\bfseries\color{accent} Five Readers on \emph{The Turn}}\\[0.5em]
{\large A referee exercise on \emph{Finding Mind}, draft 9}\\[1.2em]
{\small Spivak $\cdot$ Jaimungal $\cdot$ Goertzel $\cdot$ Tegmark $\cdot$ Bach}\\[0.4em]
{\small\color{rule} with a synthesised revision plan}\\[0.8em]
{\footnotesize\color{rule} 5 August 2026}
\end{center}

\vspace{1.2em}
\hrule
\vspace{1.2em}

\section*{How to read this}
\addcontentsline{toc}{section}{How to read this}

Five reviewers, each reading the same 570-page draft with different
commitments, each finding different soft spots. The voices are
simulated; the criticism is not. Where a reviewer speaks in the first
person, treat the idiom as a device for keeping the commitments
distinct and the objection sharp --- nothing here should be attributed
to the actual person named.

Scope is \emph{The Turn} only: Part~I (Machinery, ch.~5--17),
Part~II (Ecology as Cognitive Architecture, ch.~18--21), Part~III
(Situating the Learner, ch.~22--30), Part~IV (Why Scientists Get
Misled, ch.~31--40), Part~V (Entropy, Choice, and the Break from
Turing, ch.~41--52), Part~VI (Origins of Life, ch.~53--61), and the
Conclusion. The Pledge and The Prestige are referred to only where The
Turn depends on them.

Locations are given by source file and label, since you asked for the
\LaTeX{} source. The build tree used was
\loc{draft2}$\to\cdots\to$\loc{draft9} overlaid onto
\loc{ThePledge/}, \loc{TheTurn/}, \loc{ThePrestige/}, with
\loc{TheTurn/turn\_causality\_ch02.tex} deleted, per
\loc{INTEGRATION-NOTES-7.md}.

\vspace{0.6em}
\begin{flag}
\fnd{Five findings recur across reviewers} and are collected in
\S\ref{sec:cross} so they are not lost in the individual reports.
Three are internal contradictions the draft can repair without new
research. One is a broken proof. One is a piece of apparatus that was
retracted in one chapter and is still load-bearing in two others.
\end{flag}

\clearpage
\tableofcontents
\clearpage

\section{David Spivak reads the machinery}

\begin{voice}{The commitments}
Applied category theory. If a construction is called a functor, a
monad, a fibration, or an adjunction, the coherence data should be
exhibited or the word should be withdrawn. Universal properties are
what make a categorical framing earn its keep; without them the
diagrams are decoration.
\end{voice}

\subsection*{What works, and why I am reading the rest carefully}

Chapter~9, \emph{The Category of GSLTs}, is the best piece of
category theory in the book and I want to say so before I complain.
The three defects of the naive definition --- the constant map, the
fact that Milner's encoding is not a map of terms, the ambient-context
problem --- are diagnosed correctly, and taking contexts as primary
fixes all three at once for a reason, not by fiat. Defining morphisms
as pseudofunctors of context multicategories that are
bisimulation-preserving \emph{with the target bisimulation computed
over the image of the source's contexts} is exactly right. That last
clause is the whole game and most people writing about expressiveness
get it wrong.

So I will hold the rest of the book to that standard.

\subsection*{$\Phi$ is not a fibration}

\loc{turn\_hypercomp\_ch02.tex}, \loc{sec:fibration}:
\[
  \Phi : \mathbf{GSLT}_{\mathrm{TC}} \times \mathrm{Ord}
  \;\longrightarrow\; \mathbf{GSLT}.
\]
This is a functor \emph{out of} a product. A Grothendieck fibration is
a functor $p : \mathcal{E} \to \mathcal{B}$ with cartesian lifts. What
you have written down is an indexed family --- possibly a pseudofunctor
$\mathrm{Ord}^{\mathrm{op}} \to \mathbf{Cat}$, if you pick a direction
--- and the Grothendieck construction would produce a fibration from
it. It is not itself one.

Worse, the direction is overdetermined. You give both
$\pi_\alpha : \Phi_{\alpha+1} \to \Phi_\alpha$ and
$\iota_\alpha : \Phi_\alpha \to \Phi_{\alpha+1}$, and $\mathrm{Ord}$
is a poset, so a functor out of it can only be one of these. You have
described an adjoint pair or a section--retraction, and neither is
named. If $\iota_\alpha \dashv \pi_\alpha$ --- which I suspect is what
you want, since inclusion adds and projection forgets --- then say so
and check the triangle identities. The whole tower would then be a
tower of adjunctions, which is a much better object than what is
currently on the page and costs you nothing you do not already have.

\begin{flag}
\fnd{The book contains two objects called ``the fibration'' and they
are not the same object.} \loc{turn\_hypercomp\_ch02.tex} asserts
$\Phi$ over $\mathbf{GSLT}_{\mathrm{TC}} \times \mathrm{Ord}$.
\loc{turn\_ndo\_ch02.tex} open problem~(ii) asks someone to
``exhibit the fibration $p : \mathrm{Hyp} \to \mathbf{GSLT}$
rigorously and test whether angelic and demonic resolution are the
Lawvere adjoints $\exists \dashv (-)^{*} \dashv \forall$ to
reindexing.'' The second is the right object and is flagged as
unbuilt; the first is the wrong object and is asserted as built.
A reader who takes the chapters in order meets the assertion first.
\end{flag}

The Lawvere-adjoint question in that open problem is the correct one
and I would push harder on it than the draft does. Angelic and demonic
resolution are $\exists$ and $\forall$ over the fibres of the pairing
family; if the fibres have the requisite limits and colimits you get
the adjoint triple and with it a Beck--Chevalley condition, which is
precisely the coherence statement that ``choice strength is stable
under reindexing'' --- which is \emph{the same question} as
Question~\loc{ndo:q:crux}, the stability of rigidity under descent, and
the same question again as subject reduction in
\loc{turn\_chs\_ch01.tex}. The book says twice that two of these are
the same question. It is three, and the third framing is the one that
tells you what to prove.

\subsection*{Three monads, three adjunctions, one missing distributive law}

\loc{turn\_gslt\_ch06.tex} gives $\mathrm{Cost}$,
\loc{turn\_gslt\_ch07.tex} gives $\mathrm{Hist}$,
\loc{turn\_gslt\_ch10.tex} gives $\mathbb{K}$, and each is resolved by
a free/forgetful adjunction. Good. The unit and multiplication for
$\mathrm{Cost}$ are given explicitly and the laws are derived from the
two constituent monoids, which is the right way to do it. The
non-idempotence remark is exactly the observation a resource monad
deserves.

But the book applies these monads \emph{together} everywhere ---
Part~III decorates a cost-accounted theory that also has a history,
Part~V meters a logged theory --- and nowhere is a distributive law
$\mathrm{Cost} \circ \mathrm{Hist} \Rightarrow \mathrm{Hist} \circ
\mathrm{Cost}$ (or the reverse) stated, let alone verified. Without
one, ``a cost-accounted theory with recorded histories'' is not a
well-defined object; it is two constructions applied in an order
nobody has fixed.

And you have the material to know this matters.
\loc{turn\_entropy\_ch01.tex}, \loc{rmk:two-erasures}, says the two
counits do not factor through each other and that ``the failure to
factor is graded --- there is a lattice of intermediate laxities.'' That
remark is, read categorically, the observation that no distributive law
exists in either direction and that what you have instead is a family
of lax interchange maps. \fnd{That is a genuinely interesting
mathematical situation and the book treats it as an aside in a chapter
about entropy.} It belongs in Part~I, stated as a proposition, with
the lattice of laxities exhibited. It would also make the entropy
chapter's central move --- that the gap between the erasures is ``a
supply of positions'' --- rest on something rather than on a metaphor.

\subsection*{The compactness proof is wrong}

\loc{turn\_sciencebot\_ch06.tex},
\loc{thm:compact-consensus}. You define
$\varphi_n$ so that it is satisfied by agents that have completed
\emph{at least} $n$ oscillation steps. Then
$\llbracket\varphi_1\rrbracket \supseteq
\llbracket\varphi_2\rrbracket \supseteq \cdots$, so
$\bigcup_{n \le N} \llbracket\varphi_n\rrbracket =
\llbracket\varphi_1\rrbracket$, and an agent that has oscillated more
than $N$ times is \emph{in} that union. The claim that it is
uncovered is false, so there is no contradiction with compactness and
the theorem is unproved.

The intended argument is presumably the other way round: take the
complements, or index so that $\varphi_n$ says ``has completed
\emph{fewer than} $n$ steps,'' giving an increasing open cover of the
non-oscillating agents whose finite subcover bounds the depth. That
version may well work. As written it does not, and this is the theorem
the Conclusion singles out as ``the sharpest result of that part''
(\loc{turn\_conclusion\_ch01.tex}).

\subsection*{Scott or Stone? The book cannot have both}

Same chapter, \loc{def:log-top}. You generate the topology from HML
extensions as a subbasis and then say ``this topology is the standard
\emph{Scott topology}\,\ldots{} sober, $T_0$, and quasi-compact.''

If your HML has negation --- and \loc{turn\_causality\_ch05.tex} and the
namespace logic of \loc{ch:sci} certainly do --- then every subbasic
open is also closed, the space is Hausdorff and totally disconnected,
and you are in \emph{Stone} duality, not Scott. Points are closed, the
specialisation order is discrete, and ``sober, $T_0$'' understates it
badly. If you drop negation to get a genuine Scott topology, you lose
the ability to say $\neg\varphi$, which the consensus theorem quantifies
over.

This is not pedantry about which duality to cite. The two settings give
different answers to ``what does compactness of a population mean?''
In the Stone setting compactness is a statement about a profinite
limit and connects directly to the ultrametric structure that
\loc{ch:sci} \S\loc{sci:sec:lattice} already uses for the hypothesis
space --- which would unify two parts of the book that currently do not
speak. \fnd{Decide which topology, and the sciencebot part gets shorter
and better.}

\subsection*{Smaller things, in the order I met them}

\begin{itemize}[leftmargin=1.2em,itemsep=3pt]
\item \loc{turn\_gslt\_ch05.tex} defines \emph{hosting} as ``faithful''
      and \emph{exhausting} as ``dense'' in parentheses, then never
      uses the categorical notions again. If they are faithfulness and
      density of a (pseudo)functor, say so precisely; if they are
      analogies, drop the parentheses.
\item \loc{prop:assembly-index} (\loc{turn\_origins\_ch02.tex}) is
      labelled a proposition and has no proof. It is a definition ---
      or, if it is a claim about Assembly Theory's assembly index, it
      needs a translation functor and a preservation theorem. See
      Tegmark, \S\ref{sec:tegmark}, who is angrier about this than I am.
\item \loc{turn\_gslt\_ch10.tex}'s monad proposition asserts a
      free/forgetful adjunction between $\mathcal{L}$ and
      $\mathcal{L}_\Sigma$ and gives the action on morphisms by
      transport, but does not say what the free extension is free
      \emph{on}. For $\mathrm{Cost}$ you did say. Say it here.
\item Part~I promises ``three constructions, three monads, three
      free/forgetful resolutions'' and delivers three propositions of
      quite different strength: $\mathrm{Cost}$ is worked,
      $\mathrm{Hist}$ is sketched in four sections and 546 words, and
      $\mathbb{K}$ is asserted. A status column of the kind
      \loc{turn\_causality\_ch10.tex} uses so well would help here too.
\end{itemize}

\section{Curt Jaimungal reads for the crux}

\begin{voice}{The commitments}
Where is the load-bearing claim, stated so a smart non-specialist can
hold it and then argue with it? How does it differ from the four other
frameworks the audience already knows? And --- the question I ask
everyone --- what would have to be true for you to abandon it?
\end{voice}

\subsection*{You have three theses and the book never picks one}

Reading The Turn end to end, I can find three candidate load-bearing
claims, each stated as though it were the point.

\begin{enumerate}[leftmargin=1.4em,itemsep=3pt]
\item \emph{Mind is bounded by manufacture.}
      \loc{turn\_origins\_ch08.tex}, \loc{thm:ai-bound}, and the
      Conclusion's closing paragraphs. An ecology assembled by $n$
      joining steps resolves its world at best to $n$-step
      bisimulation.
\item \emph{Consciousness is a coordinate, not a threshold.}
      \loc{turn\_ndo\_ch02.tex} opening,
      \loc{turn\_hypercomp\_ch05.tex}. Whether every agent has a
      position in the Weihrauch lattice or crosses into a stage.
\item \emph{The noumenal survives as a budget constraint.}
      \loc{prestige\_ch01.tex} \S\loc{sec:kant-affordable} --- and you
      have said elsewhere this is the pivot of the whole book.
\end{enumerate}

The third is the one that would make an audience sit forward, and it
is in The Prestige. \fnd{The Turn builds it, states it in passing in
\loc{turn\_scope\_ch01.tex} \S``What this is for,'' and hands the
payoff to a section the reader reaches two hundred pages later.} A
reader who stops at the end of The Turn --- and many will, because The
Turn is where the mathematics is --- never learns what the framework
was for.

The Overarching Introduction opens on ``what is the minimum
mathematical structure needed to ground physics, knowledge, life, and
mind?'' That is a scope statement, not a thesis. Scope statements
promise; theses commit.

\subsection*{The comparison the book owes and does not pay}

Three sentences in \loc{turn\_hypercomp\_ch08.tex} handle Friston and
Tononi, and that chapter is visibly stale --- it still describes the
material as ``the three papers in this series'' and calls the ordinal
treatment of consciousness ``novel'' after two later chapters have
argued the ordinals are wrong. Nothing anywhere engages Wolfram's
ruliad, Levin's work on scale-free cognition, global workspace
theories, or --- most conspicuously --- Assembly Theory's own critics,
in a part that spends a hundred pages engaging Assembly Theory.

What the book needs is one page, early, in this shape:

\begin{center}
\small
\begin{tabular}{@{}p{2.5cm}p{3.6cm}p{5.6cm}@{}}
\toprule
\textbf{Framework} & \textbf{Its measure} & \textbf{Where this book disagrees} \\
\midrule
IIT & $\Phi$, integration over a partition &
  Consciousness is not integration but affordable resolution; the
  partition problem becomes the individuation fixed point of
  \loc{ch:eng} \\
FEP / active inference & Free-energy minimisation &
  Prediction error is priced in a conserved currency the agent can run
  out of; there is no scorekeeper \\
Assembly Theory & $\AI$ and $\CN$ &
  Both are supplied a region and an observer; the two stop being
  independent \\
Ruliad / Wolfram & Computational irreducibility, observer coarseness &
  The coarseness has a budget, and the budget is a physical quantity \\
Global workspace & Broadcast and access &
  Access is metered; \loc{eng:prop:fragile} prices the broadcast at
  $p^{d}$ \\
\bottomrule
\end{tabular}
\end{center}

Every row of that table is already argued somewhere in The Turn. None
of them is collected.

\subsection*{``Precisely enough to be wrong'' appears four times}

The phrase is in the Turn frontispiece, in the Overarching
Introduction, and twice more including the last line of the
Conclusion. Repetition of a defensive claim reads as anxiety, and the
book does not need it: \loc{turn\_causality\_ch10.tex}'s
\textsc{con}/\textsc{thm}/\textsc{cnj} status column does the work
honestly and is the single best editorial decision in the draft.
\fnd{Say it once, and extend the status column to every part.}

Right now the reader can audit Part~III's claims because that part has
the column, and cannot audit Part~V's or Part~VI's because they do not.
Part~V is where the boldest claims live.

\subsection*{The words \emph{sentience} and \emph{consciousness}}

\loc{turn\_hypercomp\_ch04.tex} defines sentience as a scalar running
integral of predictive accuracy that gates the budget. That is not
what anyone means by sentience; it is closer to what an economist
would call solvency and what a psychologist would call, at most, a
mood-congruent cognitive-resource effect. \loc{ch05} defines
consciousness as oracle access, which is a claim about computational
power.

You are entitled to appropriate the words. But an appropriation costs
a page of defence and currently costs a sentence. Either write the
defence --- here is what ordinary usage tracks, here is why the
technical notion deserves the name, here is what is lost --- or rename.
``Affect budget'' and ``resolution coordinate'' would cost the book
nothing except the frisson of the titles.

\subsection*{The entry problem}

570 pages, and the best material --- Part~II, the ecology --- starts at
p.~141 behind seventy pages of category theory. The Overarching
Introduction tells the reader they may begin at Part~II, but nothing
structural supports the advice: Part~II's four chapters open with
bridges that reach back into Part~I's notation, and the reader has no
signal about which Part~I constructions they will actually need.

Two cheap fixes. Give Part~II a two-page ``what you need from Part~I''
opening that states the four or five constructions it uses and where
they are. And move one worked thing forward --- the noughts-and-crosses
world of \loc{turn\_rholife\_ch02.tex} is concrete, has numbers, and
would function as an overture if a compressed version of it opened The
Turn instead of the machinery.

\subsection*{What would change your mind?}

The book answers this well in exactly one place and badly everywhere
else. \loc{turn\_entropy\_ch01.tex} \S``What this buys'' says: if every
aperture in a physically realisable system carries only finite
contention, the framework predicts no physical system exceeds the
Turing computable, ``and the question would be settled the other way.''
That is a real falsification condition and it is stated in a
subordinate clause. \fnd{It should be a section, because as Tegmark
will point out, the antecedent looks true.}

Elsewhere: what observation would refute ``mind is bounded by
manufacture''? What would refute the identification of assembly index
with minimum path length? The book's own answer --- that these are
constructions, true by definition of the objects --- is honest, and it
means the falsifiable content is concentrated in Part~VI's three
empirical remarks (\loc{scm:rmk:biosig},
\loc{scm:rmk:shed}, \loc{scm:rmk:groupsize}). Those are worth a chapter
of their own, and they are currently remarks.

\section{Ben Goertzel reads for the architecture}

\begin{voice}{The commitments}
Does this build anything? An account of mind that cannot be turned
into a running system is a philosophy of mind, which is fine, but this
book keeps promising an architecture. Show me the learning algorithm,
the inference control, the credit assignment, and the path to
open-ended novelty.
\end{voice}

\subsection*{The confinement result is the best thing here and it is
half-used}

\loc{sci:prop:confinement} --- a walk whose steps are all shorter than
$2^{-n}$ never leaves the ball of that radius --- combined with
\loc{sci:rem:nogradient} --- the hypothesis space is an ultrametric
tree and admits no gradient --- is a genuinely important negative
result about incremental learning, and I do not think it is widely
appreciated outside this book. Gradient-free, ultrametric,
provably confined: that is a precise statement of why hill-climbing in
program space does not work, and it is derived rather than observed.

The book then draws the right conclusion --- the unit of learning has
to move --- and hands the escape entirely to recombination.
\loc{turn\_rholife\_ch01.tex} \S\loc{sci:sec:repro} does this in
about two thousand words, and the crossover mechanism (the germ is a
quotation, the normal form is a tree, a locus is a context, crossover
swaps subtrees at loci) is the single most exciting construction in
the book.

\fnd{It should be a chapter. Possibly a part.} Genetic programming has
thirty years of results on exactly this operation --- destructive
crossover rates, bloat, schema theory, the ineffectiveness of
subtree crossover on typed program trees --- and none of it is cited.
More importantly, your version has something GP does not: the loci are
\emph{contexts} in a calculus with reflection, so the type discipline
that GP lacks is available for free from the OSLF construction. That
is a publishable result on its own, and it is currently a subsection.

\subsection*{Where is the inference?}

\S\loc{sci:sec:search} gives the geometry of the hypothesis space, the
constraints any revision policy must satisfy, and ``a catalogue of
strategies,'' and then explicitly declines to fix a policy
(\loc{sci:S1.4}: ``We are \emph{not} claiming to have fixed a policy,
proved a convergence theorem, or said anything about statistical
efficiency''). For a research note, defensible. For a book chapter in
a part called \emph{Ecology as Cognitive Architecture}, it is the hole
where the architecture goes.

You have the pieces for a real answer and they are scattered:

\begin{itemize}[leftmargin=1.2em,itemsep=3pt]
\item The relaxation lattice (\loc{sci:S5}) is an inference structure.
      Generalisation forgets a conjunct, specialisation adds one ---
      that is version-space learning, and version spaces have known
      complexity results.
\item The assay is a guarded receipt, so hypothesis evaluation is
      \emph{reduction}, not interpretation. That means inference
      control is scheduling, and you have an entire part
      (\loc{ch:apr}, \loc{ch:agy}, \loc{ch:chs}) about the strength of
      schedulers. Nobody has connected them. \fnd{Inference control in
      this framework is a choice-type discipline.} That is a
      one-paragraph observation that would tie Part~II to Part~V and
      neither part makes it.
\item Credit assignment is the foraging inequality
      (\loc{sci:sec:game}): a hypothesis that says where the tokens are
      is paid in tokens. That is the answer to the scorekeeper problem
      and it is genuinely elegant. But it gives you credit assignment
      only for hypotheses that were \emph{acted on}. What about the
      counterfactual --- the hypothesis that would have paid?
\end{itemize}

\subsection*{No relationship to MeTTa or Hyperon, and you built MeTTaIL}

\emph{MeTTa} appears once in the whole of The Turn, in a table in
\loc{turn\_causality\_ch10.tex}. Given that the same research
programme produced a MeTTa language-definition platform, and given
that Hyperon's Atomspace is a metagraph with rewriting and Distributed
Atomspace pattern matching, the absence is conspicuous to anyone who
knows both bodies of work.

Two things I would want said explicitly. First: is a MeTTa program a
GSLT in your sense, and if so what does the OSLF construction emit for
it? That is a concrete, checkable question and answering it would give
the framework a second worked instance besides rho, which it badly
needs --- right now every construction is illustrated on rho and a
reader cannot tell which results are about GSLTs and which are about
rho. Second: Atomspace is a shared namespace with concurrent readers
and writers, which is exactly your \loc{ch:com} setting. Your
composition-of-learners results --- disjointness is factorisation,
overlap is graded, non-interference is not stable --- are results about
Atomspace-style architectures and nobody in that community knows it.

\subsection*{The assembly-index bound is a bound on a closed system}

\loc{thm:ai-bound} says an ecology of assembly index $n$ resolves its
world at best to $n$-step bisimulation, and
\loc{rmk:power-caveat} carefully says this bounds resolution and not
reach. Good. But the slogan --- ``Mind is bounded by manufacture'' ---
overstates the theorem it summarises, in a specific way the book should
address.

Minds build instruments. A microscope is manufacture that raises the
resolution of the thing looking through it without raising the
assembly index of the looker. In your terms: an ecology can construct a
term whose interaction with the environment yields distinctions the
ecology could not make unaided, and then read the result at low depth.
That is what every scientific instrument is.

So $\AI(E)$ bounds what $E$ resolves \emph{by itself}, and minds are
distinguished precisely by not doing things by themselves. The repair
is not hard --- the bound presumably applies to $E$ together with its
instruments, which is a larger ecology with a larger index --- but it
has to be said, because as stated the theorem's most natural reading is
false and the slogan makes the false reading the memorable one.

\subsection*{Open-endedness}

Nothing in the framework generates genuinely new hypotheses. The
hypothesis language is emitted once by the OSLF construction from a
fixed presentation, and the learner then searches a lattice that was
determined before it was born. \loc{sci:S1.3} presents this as a
virtue: the language is ``generated, not designed,'' adequate for
context-labelled bisimulation and no more.

But a mind that can extend its own presentation --- adding a sort, a
rule, an equation --- is a different and more interesting object, and
your machinery does not obviously forbid it. Reflection means a term
can quote code; the germ mechanism means quoted code is heritable;
\loc{ch:eng}'s engines mean new token colours can appear. \fnd{What
would it take for a learner to extend the GSLT it lives in?} If the
answer is ``nothing, this is already available,'' say so, because it is
the difference between a framework for bounded learners and a framework
for open-ended ones. If the answer is ``the framework forbids it,''
that is a substantive limitation and should be in the open problems.

\subsection*{One small thing that bothered me}

\loc{scm:rmk:oneg}: a mind with $k$ skills does not pay the germ level
$k$ times, because ``the germ quotes the whole individual's code, and
quotation does not distribute over the parts.'' That is right, and it
is also the argument for why modular acquisition of new skills should
be cheap in this framework --- which is a claim about transfer learning
and lifelong learning that nobody makes. \loc{turn\_rholife\_ch01.tex}
\S``Why continual learning falls to the ecology'' is nearly there.
Connect them.

\section{Max Tegmark reads the physics}
\label{sec:tegmark}

\begin{voice}{The commitments}
Numbers, units, and a prediction that could come out wrong. A
correspondence between a mathematical structure and a physical one is
worth having only if it constrains something. And when a book says
four times that its claims are ``stated precisely enough to be
wrong,'' I want to know which one I should try to falsify first.
\end{voice}

\subsection*{Start with the thing I would try to falsify, because you
have it and you buried it}

\loc{turn\_origins\_ch09.tex}, \loc{scm:prop:geom} and
\loc{scm:rmk:falsepos}. An instrument whose resolution falls short of a
population's separation depth by $k$ does not miss copies --- it
reports $b^{k}$ times too many. And the inflation is largest exactly
where the population is most diverse, which is exactly where selection
has been most active. So the error is not noise; it is a bias aligned
with the signal.

\fnd{This is the best result in draft 9 and possibly in the book, and
it is a remark inside a subsection.} It is a concrete, quantitative,
practically consequential claim about the leading experimental
biosignature programme, and it is stated in a form a mass
spectrometrist could act on. Cronin's group measures $\CN$ by counting
identical fragmentation spectra; two molecules with indistinguishable
MS2 spectra are counted as copies. That is precisely a finite-resolution
equivalence, and your proposition says it inflates the count
exponentially in the shortfall.

The book then says ``we do not know how large the effect is in
mass-spectrometric practice and we would want it estimated.'' You can
estimate it. A tandem MS instrument has a known number of
distinguishable fragmentation channels; a molecular population has a
known isomer diversity for a given formula. Even an order-of-magnitude
Fermi estimate --- is $k$ typically 1, or 5? --- turns this from a
philosophical caution into a number somebody has to answer. Do the
estimate, get it wrong if necessary, and put it in the abstract of
whatever paper this becomes.

\subsection*{$10^{80}$ agents}

\loc{turn\_sciencebot\_ch05.tex}: a population of $|\mathcal{P}_0|
\approx 10^{80}$, ``comparable to the number of quarks in the
observable universe,'' with the remark that the scale is chosen so that
``every cluster of agents is itself a large enough system to implement
a complex agent.''

Two problems, and the second is the serious one.

First, $10^{80}$ is the baryon count; the quark count is about three
times that, and both are dwarfed by the photon count ($\sim 10^{89}$).
Minor, but it is the kind of thing that makes a physicist stop
trusting the rest.

Second, and fatally for the framing: if there are $10^{80}$ agents in
the observable universe, each agent has on average the mass--energy of
one proton. A proton cannot implement a complex agent. So the remark
justifying the scale --- that clusters at every level are large enough
to be agents --- is inconsistent with the scale it justifies. You
cannot have $10^{80}$ agents and have the agents be complicated.

The repair costs nothing, because you do not need the number. What the
argument needs is: \emph{a population large enough that the
communication diameter exceeds the coherence length}, whatever those
are. State it that way, scale-free, and the result gets stronger
rather than weaker --- it now applies to a bacterial colony, a research
community, and a galaxy-spanning civilisation alike, which is
presumably what you meant. Right now the physical gloss is doing
nothing but attracting fire.

\subsection*{The retracted apparatus is still load-bearing in two places}

\loc{turn\_causality\_ch09.tex} is admirable. It states the complex
decoration, states the obstruction --- bisimulation is branching-time
and refuses to pool runs, so no assignment of weights makes paths
annihilate --- explicitly retracts the density-matrix proposal, and
names three routes out without pretending to have chosen. That is how
to withdraw a claim.

The problem is that two later chapters never got the memo.

\begin{flag}
\fnd{\loc{turn\_sciencebot\_ch05.tex} defines message fidelity as
$|W(\gamma)|^{2}$, the squared modulus of the path amplitude.} The
whole of Part~IV --- communication radius, degradation at scale,
coherence clusters, the nerve, the apparent ontology --- is built on
that definition. It is the retracted apparatus.

\fnd{\loc{turn\_hypercomp\_ch05.tex} defines $\epsilon$-approximate
oracle access via $|\langle \mathtt{halt} \mid A \otimes
\lceil(M,x)\rceil\rangle|^{2}$}, ``where probability is measured by
the path integral.'' Same apparatus. And \loc{sec:hyper-gaps} Gap~7
compounds it by asking for the oracle interaction amplitude to be made
precise ``via the quantum Gillespie algorithm of Part~I'' --- which
\loc{turn\_causality\_ch09.tex} declines to call a quantum dynamics.
\end{flag}

Neither use actually needs amplitudes. Part~IV needs a fidelity that
degrades multiplicatively along a path and is bounded below by a
threshold; the stochastic instance ($\Semi = \mathbb{R}_{\ge 0}$),
which the book says is unobstructed and is what the implementation
runs, supplies exactly that. Rewriting
\loc{turn\_sciencebot\_ch05.tex} in the stochastic semiring is a
half-day of work and removes the book's largest internal inconsistency.
Part~V's graded consciousness needs a probability of correct answer,
which is the same fix.

\subsection*{``Least action is the tropical path integral''}

\loc{turn\_causality\_ch06.tex}, \loc{prop:leastaction}. I want to be
careful here because the chapter is a substantial improvement on what
it replaced --- the honest paragraph explaining that
$S = \sum \log \mathrm{dec}$ was ``a change of variable, not a
variational principle'' is exactly right and rare.

But the replacement oversells too, in a different direction. In the
tropical semiring $\oplus = \min$, so the transition weight is the
minimum over paths, and the proof is one line: ``immediate from
$\oplus = \min$ and $\otimes = +$.'' That is true and it is a
restatement of the definition of the tropical semiring. What it gives
you is \emph{Bellman's principle} --- the shortest-path structure of
dynamic programming --- which is a real and useful thing, but it is
optimal control, not Lagrangian mechanics.

The classical principle is $\delta S = 0$: the realised trajectory is a
\emph{stationary} point of the action over a continuum of nearby
trajectories, and most physically realised trajectories are saddle
points, not minima. Your construction has no continuum, no variation,
no Euler--Lagrange equation, and no $\hbar$. Calling $\min_\gamma$
``the principle of least action'' invites a physicist to look for
those four things and conclude you do not know they are missing.

The honest and still-interesting statement: \emph{in the tropical
semiring the path integral computes the Bellman value function, and
the classical variational principle would be recovered in a continuum
limit that has not been constructed.} Then put the continuum limit in
the open problems, where the Legendre transform already sits, correctly
withdrawn.

\subsection*{Landauer, without a temperature}

\loc{turn\_causality\_ch07.tex}: the erasure a computation performs is
bounded below by the holonomy of the energy form, and ``the two
thresholds coincide.'' This is presented as the most satisfying result
in the chapter and I understand why --- structurally it is lovely.

But Landauer's bound is $k_{B}T\ln 2$ per bit. It has units. It has a
temperature. Your holonomy $\theta$ is a dimensionless accumulation
around a cycle in a conversion graph. Saying the two thresholds
coincide is, as far as I can tell, the claim that two abstract lower
bounds happen to be the same abstract quantity --- which is a theorem
about your formalism, not a derivation of Landauer's principle.

To make it a physical claim you need a temperature, and the framework
has an obvious candidate it has not used: the free-energy structure of
\loc{turn\_causality\_ch07.tex} \S``The second law'' already
distinguishes total value from available value, and the ratio of
dissipation to entropy production \emph{is} a temperature in every
other thermodynamic formalism. \fnd{Define $T$ in the framework,
recover $k_{B}T\ln 2$ with the right constant, and this becomes a
result rather than a resonance.} If the constant does not come out,
that is worth knowing too.

\subsection*{Choice: two different lattices, one word}

This is my most serious technical objection and it goes to the thesis
of Part~V.

\loc{turn\_entropy\_ch01.tex} \S``Choice is what fills the gap'' argues:
picking a winner from a contending family is a choice function;
asserting a complete schedule exists is an instance of the axiom of
choice; ``finite contention is free, contention indexed by $\omega$
requires countable or dependent choice, class-sized families require
the full axiom.'' \loc{turn\_ndo\_ch01.tex} \S``A lattice, not a
tower'' then locates choice strength in the \emph{Weihrauch} lattice
of Brattka--Gherardi--Pauly.

These are different structures measuring different things.

\begin{itemize}[leftmargin=1.2em,itemsep=3pt]
\item AC-strength is set-theoretic: which fragment of choice you must
      add to ZF. It is graded by cardinality of the index family.
\item Weihrauch degree is computability-theoretic: the uniform
      reducibility of partial multivalued functions on represented
      spaces. \textsf{LLPO}, which the chapter cites as the home of
      demonic binary choice, is not an instance of AC at all --- it is
      a semi-classical logical principle, and its Weihrauch degree
      concerns the uniform computability of a selection, not the
      existence of a set.
\end{itemize}

The two do interact --- constructive reverse mathematics is exactly the
study of that interaction --- but the interaction is a research area,
not a substitution. \fnd{Part~V slides between them and the slide is
load-bearing}, because the AC framing gives you the transfinite grading
that motivates the ordinal tower, and the Weihrauch framing gives you
the lattice that supersedes it. One paragraph naming both structures
and stating which one the book uses --- and what is being conjectured
about the relation between them --- would cost half a page and would
convert an equivocation into an open problem.

\subsection*{And on the book's own terms, finite contention settles it}

\loc{turn\_entropy\_ch01.tex}, last section: ``if it turned out that
every aperture in a physically realisable system carried only finite
contention, then the framework would predict that no physical system
exceeds the Turing computable, and the question would be settled the
other way.''

Every physically realisable system carries only finite contention.
There are finitely many degrees of freedom in a bounded region with
bounded energy --- that is the holographic bound, and it is about as
well established as anything in this neighbourhood. So the antecedent
holds and the book's own falsification condition fires.

I do not think this destroys Part~V. The interesting content survives
in a different form: the \emph{location} claim --- that whatever excess
capacity exists is spent at apertures, in dissipated free energy --- is
a claim about where computational resources go, and it is true and
useful whether or not the resource is hypercomputational. What does not
survive is the framing that the ordinal tower is inhabited.

\fnd{So confront it.} Make \S``What this buys'' into a chapter that
states the finiteness objection in its strongest form --- Bekenstein,
holographic bound, finite contention --- and then says what is left.
What is left is: choice strength is a resource that is spent, its
grading is real and measurable in the Weihrauch lattice even at finite
degrees, and the tower is a mathematical idealisation whose inhabitants
are hypothesised. The book already says the last part
(\loc{turn\_entropy\_ch01.tex}: ``the tower\ldots{} remains a
construction whose inhabitants are hypothesised rather than built'').
Say it \emph{first}, not last, and Part~V becomes honest instead of
looking like it is hoping the reader will not notice.

\subsection*{The worked composite: arithmetic sound, accounting slippery}

I checked \loc{scm:sec:composite}. The numbers are internally
consistent --- $\lceil 20 \cdot 0.95^{-3}\rceil = 24$,
$\lceil 20 \cdot 0.90^{-3}\rceil = 28$, the copy series
$(1152, 48, 12, 1)$, the flat collapse $1213$, the ratios $24 : 4 :
12$, the $34.9\%$ and $57.1\%$. All correct. Two accounting problems.

\begin{enumerate}[leftmargin=1.4em,itemsep=4pt]
\item \fnd{Table~\loc{scm:tab:skills} carries two different depths in
      adjacent columns and the floor adds the other one.} The column
      $d^\ast$ (optimal \emph{integration} depth, from
      \loc{eng:prop:fragile}) reads $3,5,3,5$; the column ``separates
      at'' (modal depth at which the skill's learners become
      distinguishable) reads $4,2,3,3$. The assembly floor charges
      ``differentia: $d^\ast$ per skill, $3+5+3+5 = 16$'' and cites L4
      and L5 --- but L4's floor is $\AI \ge n$ for \emph{resolving} at
      depth $n$, which is the second column, summing to $12$. If the
      identification is intended, one sentence should say why an
      integration depth and a separation depth are the same charge.
      If it is not, the floor is $24$, not $28$, and every percentage
      in the paragraph moves.
\item \loc{scm:prop:learner} --- ``a learner costs strictly more than a
      scorer,'' $\AI(L) \ge \AI(S) + 1 + n^{-}$ --- is described as
      ``the one bound here we find genuinely interesting,'' and its
      $n^{-}$ term never appears in the only worked evaluation. Either
      it is subsumed in the $d^\ast$ charge, in which case say so, or
      the floor is missing a term.
\end{enumerate}

\subsection*{Two places to cut}

\loc{scm:rmk:groupsize} derives a critical pod size of $75$ from
$\theta = 0.2$, $\rho = 5$, an assumed isoperimetric ratio of $0.3$,
and notes that ``the literature on cohesive group size in social
mammals lives in this neighbourhood.'' Then: ``we would not want that
repeated as a result.''

If you would not want it repeated, do not print it. A number in the
Dunbar neighbourhood, derived from four invented parameters, with a
disclaimer, is exactly the artefact that gets screenshotted without the
disclaimer. The \emph{shape} of the claim --- there is a finite closure
size, it scales as $\rho/(\varepsilon\theta)$, and it is set by the
reliability of media \emph{between} individuals rather than by
anything about the individuals --- is genuinely interesting and stands
without the numeral. Keep the scaling law, delete the $75$.

Similarly \loc{turn\_hypercomp\_ch08.tex}'s ``Broader Picture'' section
should go entirely. It describes the material as three papers in a
series, credits the ordinal treatment as novel after the book has
argued it is wrong, and disposes of IIT in a sentence that says the
present approach ``avoids the measure-theoretic machinery of IIT'' ---
which reads as a boast about having less mathematics.

\section{Joscha Bach reads for the subject}
\label{sec:bach}

\begin{voice}{The commitments}
A mind is a self-organising software process that constructs a model
of its world and, crucially, a model of itself inside that world. The
interesting question about consciousness is not why there is something
it is like --- that framing smuggles in the answer --- but why a system
would \emph{represent itself} as having phenomenal states, and what
work that representation does. Show me the self-model, the attention
mechanism, and the motivational system, or you have described a very
sophisticated thermostat colony.
\end{voice}

\subsection*{You built reflection and used it for reproduction}

This is my central criticism and everything else follows from it.

The rho calculus is reflective. A term can quote itself; quotation is
inert; a name codes a strictly smaller term. \loc{turn\_rholife\_ch01.tex}
\S\loc{sci:S12.2} uses this beautifully: the germ is a quotation of the
learner's own code, transmissible and inert until dropped, and
\loc{scm:rmk:germgen} later identifies the germ with the generator of
the individual's scope. That is a genuine insight and it is doing
inheritance work.

\fnd{It is also, and the book never says this, a self-model.} A learner
that holds a quotation of its own code holds a representation of
itself, in the same language it uses for hypotheses about everything
else, in a form it can inspect without executing. The hypothesis
language is a logic of namespaces; the learner's own namespace is a
namespace; therefore the learner can \emph{form hypotheses about
itself}, priced by exactly the same assay machinery.

Nothing in The Turn does this. Search the ecology part for a hypothesis
whose subject is the learner: there is not one. The learner reasons
about its environment, forages, composes, reproduces, and never once
asks what it is. And yet:

\begin{itemize}[leftmargin=1.2em,itemsep=3pt]
\item \loc{sci:sec:exposure} --- the exposure lemma --- already prices
      \emph{being known}. That is the third-person view of the same
      quantity.
\item \loc{ch:eng} \S\loc{eng:S8} defines an individual as a fixed
      point of a partition and says explicitly that the fixed point has
      multiple solutions. A system that must \emph{choose} a solution
      to ``what am I'' is a system with a self-model, and the book
      calls this a circularity to be lived with rather than a
      cognitive function to be performed.
\item \loc{prestige\_ch01.tex} \S\loc{sec:kant-affordable} says a
      learner ``does not fail to see the distinctions it cannot afford
      --- it fails to see that they are there, and reports a world with
      fewer kinds in it than the world has.'' Apply that sentence to
      the learner's model \emph{of itself} and you have derived the
      structure of introspective illusion. The self-model is
      necessarily coarser than the self, the coarsening is not
      represented, and the system therefore reports a self that is
      simpler, more unified, and more transparent than it is.
\end{itemize}

\fnd{That is an account of why a system represents itself as having
phenomenal states, it is three sentences long, it follows from
machinery the book has already built, and it is not in the book.} It
belongs in The Turn, not The Prestige, because it is the payoff of
Part~II and Part~VI rather than a metaphysical coda.

It would also let you retire Gap~8 (\loc{sec:hyper-gaps}: ``the hard
problem\ldots{} remains open and is not claimed to be resolved here''),
which is currently an abdication dressed as modesty. You do not have to
solve the hard problem. You have to say why a budgeted reflective
learner would generate the report that produces it. You can.

\subsection*{Consciousness as oracle access does not do what the book
needs it to do}

\loc{turn\_hypercomp\_ch05.tex}, \loc{def:consciousness}: an agent is
conscious at level $\alpha$ if it has effective access to the oracle of
stage $\alpha$ and its world model reflects the finer bisimulation
quotient.

That is a definition of \emph{being cleverer}. A philosophical zombie
with a halting oracle is a zombie with a halting oracle. Nothing in the
definition makes the agent's world \emph{appear} to it, makes the
distinctions available to it \emph{as distinctions it is making}, or
distinguishes an agent that resolves finely from an instrument that
resolves finely.

The remark that follows --- ``the phenomenological richness of higher
consciousness is a direct consequence of ontological richness'' ---
asserts the bridge rather than building it. And the book knows better
elsewhere: \loc{turn\_sciencebot\_ch08.tex}, \emph{The Apparent
Ontology}, is entirely about the gap between what is there and what
appears, and it is the right chapter for this problem. Nobody connected
them.

The coordinate reframing in \loc{turn\_ndo\_ch02.tex} is a real
improvement --- ``not a threshold an agent crosses but a point it
occupies'' --- because it stops consciousness being a club with an
entrance exam. But it does not fix the underlying problem, which is
that a position in a lattice of computational strength is still a
statement about capacity, and capacity is not experience. It relocates
the assertion.

What would fix it, in your own terms: consciousness is not the
resolution an agent has, it is the agent's \emph{model of its own
resolution} --- the second-order structure by which a system represents
which distinctions it is currently making and could make. That is
buildable here: the modal depth of the formulae a learner is currently
funding is a quantity the learner could itself have a formula about,
and a system that tracks that quantity has something extremely like
attention.

\subsection*{Sentience is a budget line, not an affect}

\loc{turn\_hypercomp\_ch04.tex} makes $\varsigma$ a distinguished
component of the vectorial account, driven by a reinforcement signal
on predictive performance, gating what the agent can afford. The
remark that ``an agent that feels bad literally cannot think as well''
and the ensuing vicious cycle are the strongest part of the chapter,
and the structural parallel to depression is worth having.

But as constructed, $\varsigma$ is indistinguishable from being poor.
Four things affect is and $\varsigma$ is not:

\begin{enumerate}[leftmargin=1.4em,itemsep=3pt]
\item \emph{It is multi-dimensional.} Valence and arousal are not one
      number, and the difference matters functionally: high-arousal
      negative and low-arousal negative states produce opposite
      policies. You already have a \emph{vector} $\vec{A} =
      (A_1,\dots,A_k)$ and you singled out one component. Why not a
      projection of the whole vector, or a second-order signal about
      its rate of change?
\item \emph{It is intentional.} Affect is \emph{about} something ---
      fear of that, disgust at this. $\varsigma$ is about the agent's
      aggregate track record. The framework has the apparatus for
      aboutness: a feeling could be indexed by the namespace whose
      predictions are failing.
\item \emph{It modulates policy, not just budget.} Fear does not make
      you poorer; it makes you prefer different actions at the same
      wealth. In the book, everything affect does is done by making
      experiments unaffordable, which conflates a mood with a
      liquidity crisis.
\item \emph{It has set points and homeostasis.} $\varsigma$ has a
      threshold and a drift and no target. Motivational systems have
      regulated variables.
\end{enumerate}

And the deeper problem: \fnd{there is exactly one drive}. ``The motive
is metabolic'' (\loc{sci:S1.3}, R4) is elegant, and it solves the
scorekeeper problem honestly, which most frameworks do not. But a mind
with one drive is a bacterium. Real cognitive systems have competing,
incommensurable drives and the interesting work is in the
commensuration.

You have the machinery for this and it is in the wrong chapter.
\loc{ch:eng} gives multiple token colours, engines that convert between
them, a conversion graph whose cycles are an error-correcting code, and
--- the key piece --- the virtual token, which exists exactly when no
cycle of conversions pays for itself. \fnd{That is a theory of
commensurable drives, and the case where arbitrage is present is a
theory of incommensurable ones.} Rename nothing, add one section: the
token colours are drives, the engines are the trade-offs an animal
makes between them, and incommensurability
(\loc{turn\_rholife\_ch04.tex} \S``Incommensurability'') is what a
value conflict is.

\subsection*{No attention, and the framework is shaped like one}

A bounded learner with a tree-shaped hypothesis space, a metered
search, and a wealth-bounded branching factor
(\loc{sci:prop:branching}) is a system that \emph{must} allocate, and
\S\loc{sci:sec:search} declines to say how. The book is explicit about
this being deliberate.

I understand the reason --- you want to show which features are forced,
and policy is one of three genuinely free parameters. But the
consequence is that a part titled \emph{Ecology as Cognitive
Architecture} contains no attention mechanism, and attention is the
first thing anyone building a cognitive architecture has to supply.

What I would want, and what I think falls out: attention is the
allocation of the metered search budget across the branches of
\loc{sci:prop:branching}, and the \emph{constraints} you have derived
are already strong enough to characterise the space of admissible
attention policies even without picking one. State that as a result.
``Here is the space of possible attention mechanisms for a mortal
learner, and here is what every member of it must satisfy'' is a
stronger contribution than picking one, and it is what you already
have.

\subsection*{The unity problem, which the book almost names}

What makes the four skill-ecologies of \loc{scm:sec:composite} one
mind rather than four minds sharing a pod? \loc{ch:eng} answers:
closure. A region whose internal media are closed and whose boundary
media are not is an individual, subject to an isoperimetric condition.

That is an \emph{economic} criterion --- who can be named from outside,
who pays for what. And the book knows it is not enough:
\loc{turn\_ndo\_ch02.tex} \S``Choice and the drawing of boundaries''
distinguishes two criteria for a boundary, economic and semantic, and
says neither implies the other. \fnd{The semantic criterion is never
supplied.} So the framework has a rigorous account of what makes
something an economic individual and no account of what makes something
a single subject.

This is the unity of consciousness, and the book should name it as
such. My guess at where the answer lives, in your terms: a region is
one subject when its self-model is one model --- when the fixed-point
solution of \loc{eng:S8} that the region \emph{represents itself as}
is a single generator rather than several. That makes unity a property
of the self-model rather than of the plumbing, which is where I think
it belongs, and it makes split-brain and dissociative phenomena the
case where the plumbing and the self-model disagree. It is also
checkable in your framework, because generators are formulae and you
can ask how many there are.

\subsection*{One thing the book gets right that almost nobody does}

The claim that the true theory of an environment can be dominated by a
cruder theory that costs a fraction as much and wins nearly as often
(\loc{sci:rem:yieldbias}, and demonstrated with numbers in
\loc{turn\_rholife\_ch02.tex}) is correct, important, and rarely
stated. Most accounts of rationality treat approximation as a
regrettable concession; this treats it as the equilibrium.

That result is also, I would argue, the real answer to why minds have
the character they do --- confabulating, coarse-graining, committed to
serviceable fictions --- and it is stated as an ecological observation
rather than as a thesis about cognition. It is a thesis about
cognition. Promote it.

\clearpage
\section{Where the five disagree}
\label{sec:conflict}

A synthesis that averaged these reports would produce mush. The
disagreements are the useful part, so here they are, adjudicated.

\subsection*{Rigour versus legibility (Spivak vs.\ Jaimungal)}

Spivak wants the coherence data exhibited: distributive laws, cartesian
lifts, triangle identities. Jaimungal wants the book to be holdable by
a smart non-specialist and thinks it is already seventy pages too slow
at the front.

These are not actually opposed, because they concern different
chapters. Spivak's demands land on Part~I and the fibration; Jaimungal's
on the ordering and framing of the whole. \fnd{Resolution: do the
category theory properly and \emph{shorten} its presentation.} A
correctly stated adjunction is briefer than a hand-waved one, because
the hand-waving is where the paragraphs of apology go.

\subsection*{Buildability versus falsifiability (Goertzel vs.\ Tegmark)}

Goertzel wants an architecture that runs; Tegmark wants a number that
could come out wrong. Under a fixed writing budget these compete
directly, and the book cannot do both well in this draft.

\fnd{Resolution: Tegmark wins for The Turn, Goertzel wins for the next
paper.} The reason is that the falsifiable content already exists and
is under-presented (the overcounting bias, the depth-before-population
prediction, the $\log \CN / h$ ratio), whereas the architecture content
would be new writing of a kind that would double the ecology part. Take
the free win now. But adopt Goertzel's two cheapest asks, because they
cost paragraphs rather than chapters: state that inference control in
this framework is a choice-type discipline, and repair the
instrument-building objection to \loc{thm:ai-bound}.

\subsection*{Whether the consciousness chapters are salvageable
(Bach vs.\ everyone)}

Tegmark, Spivak and Jaimungal all treat
chapters 48--50 (\loc{turn\_hypercomp\_ch04}--\loc{ch06}) as chapters with fixable
defects: restore the coordinate revision, drop the retracted
amplitudes, decide the lattice. Bach says the defect is not fixable by
restoration, because ``consciousness is a position in a lattice of
computational strength'' is a claim about capacity and no amount of
tidying makes capacity into experience.

\fnd{Bach is right about the diagnosis and the other three are right
about the sequencing.} Restore first --- the contradiction is
indefensible and cheap to fix --- and then write the section Bach is
asking for, which is not a defence of the oracle definition but a
different claim standing beside it: that what a system reports as
phenomenal is the shortfall between its self-model and itself. That
claim is derivable from machinery already in the book and does not
require the oracle definition to be true.

\subsection*{Whether to keep the ordinal tower at all}

You decided earlier that the tower stays, as a way of introducing the
intuitions, with the lattice as the real structure. Spivak and Bach
both push back, for different reasons: Spivak because the tower as
presented is not a fibration and its repair is a different object;
Bach because a chain-shaped shadow of a lattice invites exactly the
threshold reading the book is trying to escape.

\fnd{Resolution: keep it, but demote it explicitly and early.} The
tower should be introduced as an idealisation with a stated defect ---
one sentence at first use, not a retraction three chapters later.
\loc{turn\_hypercomp\_ch02}'s missing closing section was going to do
exactly this.

\subsection*{Assembly Theory: friendly amendment or takeover?}

Tegmark reads Part~VI as a friendly amendment to Cronin and Walker
that the book undersells. Jaimungal reads it as a takeover the book has
not earned, because the answers are internal to the rho calculus and
not operational for a chemist.

Both are right about different halves. The \emph{copy-number} work is a
friendly amendment and should be presented as one, loudly. The
\emph{assembly-index} identification is a takeover and is currently
asserted as a proposition without proof. \fnd{Split the two and give
them different rhetorical registers.}

\clearpage
\section{Cross-cutting findings}
\label{sec:cross}

Six items that more than one reviewer hit, ordered by how cheap they
are to fix relative to how much damage they do.

\subsection*{1. Four of five revisions listed in
\texttt{INTEGRATION-NOTES-2.md} were never committed}

\begin{flag}
\loc{draft3/INTEGRATION-NOTES-2.md} records five revised files. Only
\loc{turn\_hypercomp\_ch07.tex} is present in the repository. Missing:

\smallskip
\begin{tabular}{@{}p{5.2cm}p{7.4cm}@{}}
\toprule
\textbf{File} & \textbf{What the notes say it contained} \\
\midrule
\loc{turn\_hypercomp\_ch01} & consciousness as a coordinate;
  organisation paragraph \\
\loc{turn\_hypercomp\_ch02} & new closing section, \emph{Two Things
  This Construction Gets Wrong} \\
\loc{turn\_hypercomp\_ch05} & the ordinal remark and the comparison
  table \\
\loc{turn\_rholife\_ch04} & forward pointer into ch.~39 \\
\bottomrule
\end{tabular}
\smallskip

Verified against the built tree: \loc{turn\_hypercomp\_ch01.tex}
contains zero occurrences of ``coordinate'' and still reads
``\emph{Consciousness} is an ordinal-valued threshold property'';
\loc{turn\_hypercomp\_ch05.tex} still has
\loc{remark}[Consciousness is Ordinal-Valued] and the table row
``Changes: By threshold (have it or not)'';
\loc{turn\_rholife\_ch04.tex} contains zero references to
\loc{ch:apr} or \loc{ch:agy}.
\end{flag}

The visible damage is a forward reference that lies.
\loc{turn\_ndo\_ch02.tex} tells the reader that consciousness in the
chapters that follow is ``not a threshold an agent crosses but a point
it occupies, with the ordinal tower\ldots{} the chain-shaped shadow of
the lattice that point really lives in.'' Two chapters later the book
says consciousness is discrete, ordinal, and changes ``by threshold
(have it or not).'' \loc{turn\_conclusion\_ch01.tex} sides with the
older view (``the first continuous, the second not'').

This is the highest-value repair in the document and it is a
restoration, not new work.

\subsection*{2. Retracted apparatus still load-bearing in two chapters}

\loc{turn\_causality\_ch09.tex} withdraws the complex-amplitude
reading. Still using it:

\begin{itemize}[leftmargin=1.2em,itemsep=2pt]
\item \loc{turn\_sciencebot\_ch05.tex}, \emph{Message Fidelity}
      $= |W(\gamma)|^{2}$ --- the foundation of all of Part~IV.
\item \loc{turn\_hypercomp\_ch05.tex}, $\epsilon$-approximate oracle
      access via the path integral.
\item \loc{sec:hyper-gaps} Gap~7, which asks for the oracle amplitude
      to be made precise via the quantum Gillespie algorithm.
\end{itemize}

Neither use needs amplitudes; the stochastic instance supplies a
multiplicatively degrading fidelity and a probability, and the book
says that instance is unobstructed and is what the implementation runs.

\subsection*{3. \texttt{thm:compact-consensus} is unproved as written}

\loc{turn\_sciencebot\_ch06.tex}. The family
$\{\llbracket\varphi_n\rrbracket\}$ is decreasing, so
$\bigcup_{n\le N}\llbracket\varphi_n\rrbracket =
\llbracket\varphi_1\rrbracket$ and the agents claimed to be uncovered
are covered. The Conclusion calls this the sharpest result of its part.

\subsection*{4. Two objects named ``the fibration''}

$\Phi : \mathbf{GSLT}_{\mathrm{TC}} \times \mathrm{Ord} \to
\mathbf{GSLT}$ (\loc{sec:fibration}) is an indexed family with
contradictory variance, asserted as built.
$p : \mathrm{Hyp} \to \mathbf{GSLT}$ (\loc{turn\_ndo\_ch02.tex}, open
problem~(ii)) is the right object, flagged as unbuilt. The reader meets
the assertion first.

\subsection*{5. \texttt{turn\_hypercomp\_ch08.tex} is stale}

Describes the material as ``the three papers in this series''; credits
the ordinal treatment of consciousness as ``novel'' after two later
chapters argue it is wrong; disposes of Friston in three sentences and
Tononi in two.

\subsection*{6. The falsifiable content is real and is filed as remarks}

Every reviewer noticed some version of this. The claims that could come
out wrong are \loc{scm:prop:geom} with \loc{scm:rmk:falsepos} (coarse
instruments inflate copy number exponentially, biased toward the
biosignature), \loc{scm:rmk:biosig} (the ratio $\log \CN / h$ should
cluster near $\log(1/p)$ for architectures in use), and
\loc{scm:rmk:shed} (lineages under scarcity shed depth before
population). Three remarks, in one chapter, near the end of a
570-page book that says four times that it is stated precisely enough
to be wrong.

\clearpage
\section{Revision plan}

Four tiers. Tier~0 is restoration and repair of things the book already
believes; Tier~1 is structural moves that change what the reader
encounters and in what order; Tier~2 is new writing that the reviewers
converged on wanting; Tier~3 is what to leave alone.

Effort is in units of a working session. ``Files'' names what changes.

\subsection*{Tier 0 --- Repairs. Do these first; none needs new research.}

\begin{center}
\small
\begin{tabular}{@{}p{0.5cm}p{6.4cm}p{4.0cm}p{1.4cm}@{}}
\toprule
& \textbf{Action} & \textbf{Files} & \textbf{Effort} \\
\midrule
R1 & Restore the four lost \loc{INTEGRATION-NOTES-2} revisions: the
     coordinate framing in ch01, \emph{Two Things This Construction
     Gets Wrong} in ch02, the ordinal remark and comparison table in
     ch05, the forward pointer in rholife ch04. Then fix the
     Conclusion's ``the first continuous, the second not''.
   & \loc{turn\_hypercomp\_ch01/02/05}, \loc{turn\_rholife\_ch04},
     \loc{turn\_conclusion\_ch01}
   & 1--2 \\
\addlinespace
R2 & Rewrite message fidelity in the stochastic semiring; rewrite
     $\epsilon$-approximate oracle access as a probability rather than
     a squared amplitude; delete or restate Gap~7.
   & \loc{turn\_sciencebot\_ch05}, \loc{turn\_hypercomp\_ch05},
     \loc{turn\_hypercomp\_ch07}
   & 1 \\
\addlinespace
R3 & Repair or withdraw \loc{thm:compact-consensus}. The likely fix is
     to index the formulae by ``fewer than $n$'' and take an increasing
     cover. If it does not go through, demote to a conjecture and
     amend the Conclusion, which currently calls it the part's
     sharpest result.
   & \loc{turn\_sciencebot\_ch06}, \loc{turn\_conclusion\_ch01}
   & 1 \\
\addlinespace
R4 & Decide Scott or Stone. With negation in HML the space is Stone;
     say so, and connect to the ultrametric hypothesis space of
     \loc{sci:sec:lattice}.
   & \loc{turn\_sciencebot\_ch06}
   & 0.5 \\
\addlinespace
R5 & Downgrade $\Phi$ from ``fibration'' to ``indexed family'', name
     the intended variance, and add a forward pointer to
     \loc{turn\_ndo\_ch02} open problem~(ii) as the object that would
     replace it.
   & \loc{turn\_hypercomp\_ch02}
   & 0.5 \\
\addlinespace
R6 & Rewrite \loc{turn\_hypercomp\_ch08} from scratch, or delete it and
     fold one paragraph into \loc{turn\_conclusion\_ch03}.
   & \loc{turn\_hypercomp\_ch08}
   & 0.5 \\
\addlinespace
R7 & Resolve the $d^{\ast}$ / ``separates at'' ambiguity in the worked
     composite; either justify the identification in one sentence or
     recompute the floor at 24 and update the two percentages.
   & \loc{turn\_origins\_ch09}
   & 0.5 \\
\addlinespace
R8 & Delete the numeral 75 from \loc{scm:rmk:groupsize}, keep the
     scaling law. Fix ``quarks'' and drop $10^{80}$ in favour of a
     scale-free statement of the diameter/coherence condition.
   & \loc{turn\_origins\_ch09}, \loc{turn\_sciencebot\_ch05}
   & 0.5 \\
\bottomrule
\end{tabular}
\end{center}

\subsection*{Tier 1 --- Structural. These change the reader's path.}

\paragraph{S1. Promote the falsifiable content into its own chapter.}
Take \loc{scm:prop:geom}, \loc{scm:rmk:falsepos}, \loc{scm:rmk:biosig}
and \loc{scm:rmk:shed} out of \loc{turn\_origins\_ch09} and make a
short chapter --- \emph{What Would Show This Is Wrong} --- at the end of
Part~VI, with the overcounting result as its lead. Add the Fermi
estimate of $k$ for tandem MS. This is the single highest-leverage
editorial move available and it costs about two sessions.

\paragraph{S2. Extend the status column to every part.}
\loc{turn\_causality\_ch10}'s \textsc{con}/\textsc{thm}/\textsc{cnj}
table is the best editorial device in the book and appears in exactly
one part. Part~V and Part~VI need it more than Part~III does, because
their claims are bolder. Then delete three of the four occurrences of
``precisely enough to be wrong'' --- the column does that work now.

\paragraph{S3. Give Part~II a two-page on-ramp.}
``What you need from Part~I'': the four or five constructions the
ecology chapters actually use, with pointers. This cashes the
Overarching Introduction's promise that a reader may begin at Part~II,
which nothing currently supports.

\paragraph{S4. Move the finiteness confrontation to the front of
Part~V.} \loc{turn\_entropy\_ch01}'s closing admission --- that if every
physical aperture carries only finite contention the framework predicts
no physical hypercomputation --- should open the part, stated in its
strongest form, with what survives stated immediately after. What
survives is the location claim, which is the part's real contribution
and is independent of whether the tower is inhabited.

\paragraph{S5. Promote crossover.}
\loc{sci:sec:repro} is a subsection carrying the book's most novel
mechanism: crossover on quoted code with contexts as loci and a type
discipline available from OSLF. Make it a chapter of Part~II. Cite the
genetic-programming literature it is implicitly contradicting.

\paragraph{S6. Split the two Assembly Theory claims.}
Present the copy-number work as a friendly amendment (it is, and it is
good). Downgrade \loc{prop:assembly-index} from a proposition to a
proposed identification, and state the missing object --- a functor from
chemical reaction systems to ciGSLTs preserving minimum path length,
with AT's reuse discipline --- as the central open problem of Part~VI.

\subsection*{Tier 2 --- New writing the reviewers converged on}

\paragraph{N1. The self-model section.} \emph{(Bach; endorsed in effect
by Jaimungal, who wants the pivot inside The Turn.)} The germ is a
quotation of the learner's own code; the hypothesis language is a logic
of namespaces; the learner's namespace is a namespace. Therefore a
learner can form hypotheses about itself, priced by the same assay. Its
self-model is necessarily coarser than itself, the coarsening is not
represented, and it reports a self simpler and more unified than it is.
Place at the end of Part~II or as the closing section of \loc{ch:eng}.
Retire Gap~8 by replacing it with this.

\paragraph{N2. Drives, in the engine chapter.} \emph{(Bach.)} Token
colours are drives; engines are the trade-offs; the virtual token is
what makes drives commensurable; arbitrage is what makes them
incommensurable. One section, no new mathematics, and it turns
\loc{ch:eng} from an economics chapter into a motivational-systems
chapter.

\paragraph{N3. The two-lattices paragraph.} \emph{(Tegmark.)} Name
AC-strength and Weihrauch degree as different structures, say which one
Part~V uses, and state the relation between them as an open problem
rather than an identification. Half a page in \loc{turn\_entropy\_ch01}
or \loc{turn\_ndo\_ch01}.

\paragraph{N4. The comparison table.} \emph{(Jaimungal.)} One page,
early in The Turn: IIT, FEP, Assembly Theory, ruliad, global workspace,
and where this book disagrees with each. Every row is already argued
somewhere; none is collected.

\paragraph{N5. The interchange lattice.} \emph{(Spivak.)}
\loc{rmk:two-erasures} says the cost and history counits do not factor
through each other and that the failure is graded. Stated properly in
Part~I, that is the observation that no distributive law exists in
either direction and that what exists instead is a lattice of lax
interchange maps. It also makes the entropy chapter's ``supply of
positions'' rest on a proposition.

\paragraph{N6. Instruments.} \emph{(Goertzel.)} One paragraph
qualifying \loc{thm:ai-bound}: the bound is on an ecology resolving
unaided, and instrument-building is precisely how minds evade it. The
bound then applies to ecology-plus-instruments, which is a larger
object with a larger index. Without this the slogan's natural reading
is false.

\paragraph{N7. Inference control is a choice-type discipline.}
\emph{(Goertzel.)} One paragraph connecting \loc{sci:sec:search} to
\loc{ch:chs}. Ties Part~II to Part~V across a gap neither part
currently bridges.

\subsection*{Tier 3 --- Leave alone}

\begin{itemize}[leftmargin=1.2em,itemsep=3pt]
\item \loc{turn\_gslt\_ch05}, \emph{The Category of GSLTs}. Spivak's
      only unqualified praise. Do not touch it.
\item \loc{turn\_causality\_ch09}. The retraction is a model of how to
      withdraw a claim. The problem is downstream, not here.
\item \loc{turn\_rholife\_ch02}, the worked game. It is the only place
      a claim is checked against a number and the numbers hold up.
\item The decision to keep the ordinal tower as intuition-scaffolding.
      Two reviewers push back; the answer is to demote it explicitly
      (R1, R5), not to remove it.
\item The Pledge. Out of scope, and left untouched by design.
\end{itemize}

\subsection*{Suggested order}

Tier~0 entire, in one pass, since R1--R3 are the ones that make the
book self-contradictory and they are all restorations or local repairs.
Then S1 and S2, which are the cheapest way to change how the book reads
to a sceptic. Then N1, which is the one piece of new writing every
reviewer wanted in some form and which puts the pivot inside The Turn
where Jaimungal argues it belongs. Everything else can wait for the
next draft.

\vspace{2em}
\hrule
\vspace{0.8em}
{\footnotesize\color{rule}
The five voices are simulated readings, constructed to keep distinct
sets of commitments in play. Attribute the criticism to the exercise,
not to the people named.}

\end{document}
